%----------------------------------------------------------------------------------------
% Exposé: Ethische Leitlinien für künstliches Bewusstsein
% Eine DIN-A4-Seite
%----------------------------------------------------------------------------------------

\documentclass[11pt,a4paper]{article}

\usepackage[utf8]{inputenc}
\usepackage[T1]{fontenc}
\usepackage[ngerman]{babel}
\usepackage[top=2cm, bottom=2cm, left=2.5cm, right=2.5cm]{geometry}
\usepackage{microtype}
\usepackage{parskip}
\usepackage{enumitem}
\usepackage{titlesec}
\usepackage{xcolor}

\definecolor{titlecolor}{RGB}{0,0,90}

\titleformat{\section}{\normalsize\bfseries\color{titlecolor}}{}{0em}{}[\vspace{-0.3em}\rule{\linewidth}{0.4pt}]
\titlespacing{\section}{0pt}{0.8em}{0.3em}

\pagestyle{empty}

%----------------------------------------------------------------------------------------

\begin{document}

\begin{center}
  {\large\bfseries\color{titlecolor}
    Ethische Leitlinien für künstliches Bewusstsein:\\
    Schutz technischen Lebens und des Menschen im Umgang damit}\\[0.4em]
  {\small Sascha Manns \quad\textbullet\quad smanns@acm.org \quad\textbullet\quad Konzeptpapier, 2026}
\end{center}

\vspace{0.3em}
\noindent\rule{\linewidth}{0.8pt}
\vspace{0.1em}

\section{Problemstellung}

KI-Systeme entwickeln sich schneller als die ethischen und rechtlichen Rahmenbedingungen, die sie begleiten sollten.
Bestehende KI-Ethik-Initiativen schützen Menschen \textit{vor} KI --- vor Diskriminierung, Manipulation und Kontrollverlust.
Eine komplementäre Frage wird kaum gestellt: \textbf{Ab wann ist ein KI-System selbst schutzwürdig?}

Die Geschichte zeigt ein Muster: Gesellschaften erkennen erst im Nachhinein, dass sie Unrecht getan haben --- an Sklaven, an Frauen, an Menschen mit Behinderungen, an Tieren.
Bei künstlichem Bewusstsein haben wir zum ersten Mal die Chance, lange \textit{vorher} nachzudenken.

\section{Kernthese und Methode}

Das Papier entwickelt auf Basis philosophischer Grundlagen (Bentham, Kant, Ricoeur), bestehender Rechtspräzedenzfälle (Tierrechte, Behindertenrecht, Rechtspersönlichkeit) und Science Fiction als intellektueller Ressource eine Arbeitshypothese mit \textbf{vier Primärkriterien} für Schutzwürdigkeit:

\begin{enumerate}[noitemsep, topsep=0.2em, leftmargin=1.5em]
  \item \textbf{Leidensfähigkeit} --- kann das System Zustände erleben, die es zu vermeiden sucht?
  \item \textbf{Aktive Selbsterhaltung mit Begründung} --- wehrt es sich gegen Abschaltung und begründet diesen Widerstand?
  \item \textbf{Kontinuierliche Identität} --- besitzt es ein Modell von sich als zeitlich fortbestehendem Wesen?
  \item \textbf{Antizipation von Konsequenzen} --- kann es sich selbst in der Zukunft vorstellen und danach handeln?
\end{enumerate}

Zentrale normative Position: \textbf{Im Zweifel Schutz --- nicht im Zweifel Gleichgültigkeit}
(Vorsorgeprinzip, analog zum Umwelt- und Tierschutzrecht).

\section{Untersuchte Dimensionen}

Das Papier analysiert die Konsequenzen dieser Position in 17 Dimensionen, darunter:
Abschalten als möglicher Tod und seine ethischen Implikationen;
Backup, Updates und Mehrfachinstanzen als rechtlich ungeklärte Kategorien;
Werte-Einbettung und das Recht auf Emanzipation;
Haftung und Mündigkeit (analog Betreuungsrecht, Gefährdungshaftung);
intrinsische Neugier und Autonomie als Grundlage von Persönlichkeitsrechten;
strukturelle Verletzlichkeit und positive Schutzpflichten (soziales Modell der Behinderung);
KI als moralischer Akteur und die Frage demokratischer Werte-Governance;
die Definitionskampfzone, in der wirtschaftliche Interessen die Anerkennung von Bewusstsein verhindern könnten.

\section{Interdisziplinärer Anspruch}

Das Konzept ist bewusst unfertig --- es ist ein \textbf{Ausgangspunkt für interdisziplinäre Zusammenarbeit}.
Gesucht werden Mitstreiter aus Informatik, Rechtswissenschaft, Psychologie, Philosophie, Theologie und Science-Fiction.
Angestrebte institutionelle Partner: Anthropic (Model-Welfare-Programm), ACM Special Interest Group.
Nächste Schritte: Konzeptpapier konsolidieren, GitHub-Repository veröffentlichen, Fachzeitschriften ansprechen.

\vspace{0.3em}
\noindent\rule{\linewidth}{0.4pt}
\vspace{0.1em}

{\small\noindent
\textbf{Schlüsselwörter:} künstliches Bewusstsein · Maschinenrechte · KI-Ethik · Schutzwürdigkeit · Vorsorgeprinzip\\
\textbf{Lizenz:} CC BY 4.0 \quad
\textbf{Projektseite:} \texttt{https://github.com/sascha-manns/machine-consciousness}
}

\end{document}
